


something about proteins in general, how they're formed, what they do, etc.

make introduction about protein modeling, aboud making models of proteins

sequence-structure-function paradigm, but substitute it with sequence-structure-interaction

much of protein science revolves arund making models that relate these different components with
each other
--> elaborate on this further

section on the binding process, with equations and whatnot. talk about the dynamic nature of it, but
also the
thermodynamic description, which averages over the dynamics


mention molecular dynamics as an attempt of a first-principles approach to study protein function 
--> first-principles (or at least close to them) because it does not hide the dynamic nature of
proteins
--> cite DESRES papers and maybe old media coverage
--> DESRES paper where they used "fraction of time ligand is bound to protein within MD trajectory"
as a proxy for
binding affinity --> example of going about it via first principles


"Molecular dynamics force fields
A force field like AMBER, CHARMM, or OPLS is built differently. It's a parametrized potential energy
function — bonds, angles, dihedrals, Lennard-Jones, Coulomb — where the parameters are fit to
reproduce quantum chemistry calculations on small molecules and experimental thermodynamic data
(densities, heats of vaporization, hydration free energies). Every atom is explicit, water is
usually explicit too (TIP3P, SPC/E, etc.), and the function returns energies in real kcal/mol. The
function is differentiable in Cartesian coordinates, so you integrate Newton's equations forward in
time with femtosecond time steps to generate a trajectory. The Boltzmann distribution at the
simulation temperature is sampled (in principle, given enough time), so you can compute
thermodynamic averages, kinetics, and free energy differences with proper statistical mechanics."


Rosetta's energy function (Alford et al. in 2017) is a weighted sum of terms that score a static
protein conformation.
The terms include both approximations of physical interactions (e.g. Lennard-Jones potentials,
hydrogen-bonding components)
and statistical regularities extracted from protein structures deposited in the Protein Data Bank
(PDB).
An implicit model model of solvent is also present, thus avoiding the need of explicitly placing
water model molecules in the structure for scoring.

"Rosetta is a scoring function: a tool engineered to rank conformations and sequences such that
native-like states score well. It's not trying to give you correct dynamics; it's trying to give
you a landscape where the global minimum is the right answer for design and prediction tasks."
--> I really like the last sentence


machine learning effectively reduces to using only statistical regularities for modeling.
--> AF2, and then Sergey's paper claiming AF2 (using MSA!) has learned an energy landscape! no need for explicit physics??
of course, the truth is more nuanced than this. in fact, in designing a machine learning solution,
one can choose exactly what information to include in the "statistical mining" performed by the model.

While it is accepted that the amino-acid sequence contains enough information to determine the lowest-energy
conformation of a protein, it is unclear how many sequence-structure pairs are required to learn such
a mapping robustly. If the available pairs are insufficient, an alternative is to include other kinds of
data that contain information about the underlying process that is to be modeled.
This is exactly what AlphaFold does: instead of predicting structure from the amino-acid sequence alone,
the model incorporates a Multiple Sequence Alignment of the seuqnce against all known other proteins.
This changes the exact function learned by the model, such that the model requires less input-output pairs to learn
a function that better represents the underlying process.


a machine learning model would be able to learn such mapping only in the limit
of high data.  

--> zero-shot!!!
--> still, it is important to attempt an understanding of the physical rules encoded by the machine
learning model, as they often reveal "shortcuts" that are not physical and are at the core of
the model's mistakes.


--> machine learning!! mention alphafold as a revelatory moment for the field

machine learning entails learning statistical patterns between paired entities

machine learning entails learning a mapping between two domains (?)

it's indeed not an accident that machine learning has been linked with statistical mechanics, as one
can view machine
learning as providing a data-driven mean-field theopry of a process 


"A mean-field theory replaces interactions with an effective field that captures the average
behavior of everything else.
A trained ML model, in a real sense, also replaces a complicated joint distribution with a tractable
surrogate
— often a factorized or low-rank one."


maybe a section on zero-shot learning as a concept? mention our perspective article on Sergey's
paper,
which is a seminal paper trying to understand the use of confidence metrics of structure-prediction
algorithms
--> notion of predicting the amount with which a design is "in-distribution"


Introduce equivariance simply as a way to simply restrict the function class learned by a model such that
it respectes known symmetry of a problem.

Introduce locality as an assumption that simplifies the mapping, while introducing approximation error
that we know to be limited given our knowledge of the underlying physics that out model is approximating.
Furthermore, insofar as the models we create are tools for design - i.e. to sample the viable protein space
in search of \textit{at least one example} that works - instead of claiming to accurately model all of
protein space, then the assumption of locality is a specification of the tool - akin to a function's
contract in programming - which does limit its use cases, but also an improved 

the paragraph above makes me think: what if I wrote a paragraph on design vs. modeling the entire landscape?



let's think about figures:
1. 


